\section{A transfer family that cannot violate the rank-distance scale}
\label{sec:four-high-connectivity-path-transfers}

The proposed general bound
\[
 P_{uv}(N)-2^{b-1}
 \le \operatorname{poly}(b)\,2^{(b+\operatorname{dist}(u,v))/2}
\]
remains open in these notes.
Repeated triangle expansion is one natural fixed-width test family.
Its exact recurrence gives a faster-decaying excess, rather than a counterexample.

\begin{proposition}[Repeated expansion beside a fixed edge]
\label{prop:four-triangle-path-transfer}
Let $N$ be a loopless normalized two-pole of rank $b$, with terminals $u,v$.
Let $x$ be a nonterminal vertex and let $f$ be an incident edge.
Partition its terminal paths into three classes, of sizes
\[
 Q=\#\{P:f\in P\},\qquad
 Z=\#\{P:x\in P,\ f\notin P\},\qquad
 A=\#\{P:x\notin P\}.
\]
Replace $x$ by a triangle, attaching its three former edges to distinct triangle vertices.
Repeat at the new vertex incident with $f$.
After $j$ expansions, the rank is $b+j$ and the path count is exactly
\[
 P_j=2^jQ+A+(j+1)Z.
\]
If $f$ is not a bridge, then $Q\le2^{b-1}$.
If this inequality is strict, the normalized path count is eventually below one half.
If equality holds, its excess is
\[
 \frac{P_j}{2^{b+j}}-\frac12
 =\frac{A+(j+1)Z}{2^{b+j}}.
\]
\end{proposition}

\begin{proof}
A simple terminal path enters a replacement triangle through exactly two external edges, or avoids it entirely.
There are only three external incidences, and neither terminal is inside the triangle.
Thus no path can have two separate segments inside the triangle.
Each path using the old vertex has two possible internal arcs after expansion.

Write $Q_j,Z_j,A_j$ for the three counts relative to the current vertex incident with $f$.
Paths using $f$ have two lifts.
A path using the other two incidences has one lift through the distinguished vertex and one avoiding it.
Therefore
\[
 Q_{j+1}=2Q_j,\qquad Z_{j+1}=Z_j,\qquad A_{j+1}=A_j+Z_j.
\]
Starting from $Q_0=Q,Z_0=Z,A_0=A$ gives the displayed formula.
Each expansion adds two vertices and three edges, increasing rank by one.

If $f$ is not a bridge, its coordinate is nonconstant on the cycle space.
Exactly $2^{b-1}$ terminal joins use $f$.
Since each terminal path is a join, $Q\le2^{b-1}$.
If the inequality is strict, the exponential deficit $2^j(2^{b-1}-Q)$ eventually exceeds $A+(j+1)Z$.
The equality case follows by subtraction.
\end{proof}

If $A+Z>0$, these graphs have a terminal path of bounded length throughout the expansion.
An original path avoiding $x$ remains unchanged.
If only $Z$ is positive, choose its lift that avoids the distinguished vertex after the first expansion.
That lift remains unchanged thereafter.
Thus even at bounded terminal distance, this critical transfer has excess only $O(j2^{-j})$ after normalization.

The equality case has a useful four-terminal interpretation.
Suppose $f=ab$ has endpoints distinct from $u,v$ and is not a bridge.
Then $K=N-f$ is a connected normalized four-pole of rank $b-1$.
The equality $Q=2^{b-1}$ holds precisely when every $\{u,v,a,b\}$-join in $K$ is a forest and none has pairing $uv\mid ab$.
Indeed, adding $ab$ turns either other pairing into one terminal path.
The excluded pairing instead gives a terminal path and a separate cycle.
Any existing cycle component also prevents the result from being a single path.

This identifies a narrower structural problem for future transfer arguments.
It does not classify such saturated four-poles or bound their residual terminal-path count.
The existing Horak--Kim construction already requires a linear polynomial factor at the half exponent.
The triangle recurrence supplies no stronger obstruction.
